\chapter{Finitude projective et sélecteur harmonique dodécaphasique}
\label{chap:finitud-anisotropia}
\index{contre-terme}\index{régulateur externe}
\index{anisotropie}\index{CMB}\index{dodécaphase}

\section{Raffinement interne du sélecteur harmonique indépendant de contre-termes externes}
\label{sec:refinamiento-sin-contraterminos}

L'affirmation selon laquelle HMT ne requiert pas de renormalisation fondamentale distingue
trois opérations : la normalisation d'états ou de probabilités, le raffinement
discret avec changement d'échelle et la renormalisation perturbative au moyen d'un
régulateur, de paramètres nus divergents et de contre-termes dépendant de la
coupure. Les deux premières sont compatibles avec HMT ; la troisième n'est pas introduite
comme primitive.

Soient
\[
 p_{m+1,m}:X_{m+1}\longrightarrow X_m
\]
les projections surjectives d'un système inverse d'espaces finis et
soient
\[
 \mathcal O_m=C(X_m)
\]
les algèbres involutives d'observables complexes. Les morphismes de
contravariance
\[
 p_{m+1,m}^{\ast}:\mathcal O_m\longrightarrow\mathcal O_{m+1},
 \qquad
 p_{m+1,m}^{\ast}f=f\circ p_{m+1,m},
\]
sont des \(\ast\)-homomorphismes unitaires cohérents. Un observable cylindrique est
une classe compatible sous ces morphismes.

Pour des mesures de probabilité, la cohérence exacte exige
\[
 (p_{m+1,m})_\ast\mu_{m+1}=\mu_m.
\]
Pour des états algébriques positifs et normalisés
\(\omega_m:\mathcal O_m\to\CC\), la condition est
\[
 \omega_{m+1}\circ p_{m+1,m}^{\ast}=\omega_m.
\]
Lorsque \(\omega_m(f)=\int f\,d\mu_m\), les deux formulations sont équivalentes.

\begin{theorem}[Compatibilité projective et défauts sommables]
\label{thm:compatibilidad-proyectiva-defectos}
Soit \(f\in\mathcal O_m\) un observable cylindrique et supposons que sa famille
de représentants soit uniformément bornée.
\begin{enumerate}
 \item Si les mesures ou les états sont projectivement cohérents, leurs
 espérances sont compatibles dans tous les raffinements.
 \item Si
 \[
 \left\|(p_{m+1,m})_\ast\mu_{m+1}-\mu_m\right\|_{\rm TV}
 \leq\varepsilon_m,\qquad
 \sum_m\varepsilon_m<\infty,
 \]
 alors les espérances forment une famille de Cauchy. La même conclusion
 vaut pour les états lorsque le défaut est mesuré dans la norme duale de
 \(\mathcal O_m^\ast\).
\end{enumerate}
\end{theorem}

\begin{proof}
Dans le cas exact,
\[
 \int_{X_{m+1}}f\circ p_{m+1,m}\,d\mu_{m+1}
 =\int_{X_m}f\,d\mu_m,
\]
et l'identité analogue pour les états est leur propre condition de cohérence.
Dans le cas approché, la dualité entre la norme uniforme et la variation
totale donne
\[
 \left|
 \int f\circ p_{m+1,m}\,d\mu_{m+1}
 -\int f\,d\mu_m
 \right|
 \leq \|f\|_\infty\varepsilon_m.
\]
La bornitude uniforme et la sommabilité des défauts transforment la somme des
incréments en une série convergente ; par conséquent, la famille est de Cauchy. La
preuve dans la norme duale est identique.
\end{proof}

Pour les actions et les mesures de référence, la condition correcte est la
cohérence des mesures de Gibbs :
\[
 (p_{m+1,m})_\ast
 \left[
 Z_{m+1}^{-1}e^{-S_{m+1}}\,d\nu_{m+1}
 \right]
 =
 Z_m^{-1}e^{-S_m}\,d\nu_m.
\]
Une simple égalité entre actions ne suffit pas, car une application de comparaison
pourrait dissimuler une soustraction ou un changement d'échelle dépendant de la coupure.

La conclusion publiable est donc restreinte à ses hypothèses :
là où existe une famille projectivement cohérente d'états ou de mesures,
les espérances des observables cylindriques uniformément bornés sont
compatibles ; avec des défauts sommables en variation totale ou en norme duale, elles forment
une famille de Cauchy. Aucun régulateur externe ni contre-terme
externe n'est requis, bien qu'il puisse exister des flux effectifs internes finis.

Dans le système projectif déclaré, la compatibilité
des espérances et le critère de Cauchy énoncés sont satisfaits exactement. Une théorie quantique des
champs complète requiert en outre la
convergence des corrélateurs, la positivité, l'unitarité, l'indépendance
cofinale et l'absence de dépendance à la coupure dissimulée dans les applications, les actions,
les mesures ou les paramètres. Un secteur qui ne converge qu'après introduction de paramètres
dépendant de la coupure réfute la version forte pour ce secteur.

\section{Sélecteur harmonique du cycle dodécaphasique}
\label{sec:selector-armonico-c12}

Soit \(\mathsf C_{12}^{\rm cyc}\) le cycle dodécaphasique local,
\[
 \phi_k=\phi_0+\frac{2\pi k}{12},
 \qquad k=0,\ldots,11,
\]
et considérons la famille de poids
\[
 w_k=\frac1{12}
 \left[1+2\varepsilon\cos(3\phi_k-\delta)\right],
 \qquad |\varepsilon|\leq\frac12.
\]
Alors
\[
 w_k\geq\frac{1-2|\varepsilon|}{12}\geq0,
 \qquad
 \sum_{k=0}^{11}w_k=1.
\]

\begin{theorem}[Support de Fourier dodécaphasique]
\label{thm:soporte-fourier-dodecafasico}
La transformée discrète
\[
 \widehat w_m=\sum_{k=0}^{11}w_ke^{-im\phi_k}
\]
s'annule sauf pour
\[
 m\equiv0,\ \pm3\pmod{12}.
\]
La phase d'origine étant absorbée dans \(\delta\),
\[
 \widehat w_0=1,\qquad
 \widehat w_3=\varepsilon e^{-i\delta},\qquad
 \widehat w_{-3}=\varepsilon e^{i\delta}.
\]
\end{theorem}

\begin{proof}
La transformée de la partie uniforme est le caractère trivial du cycle. En
écrivant
\[
 2\cos(3\phi_k-\delta)
 =e^{i(3\phi_k-\delta)}+e^{-i(3\phi_k-\delta)},
\]
les sommes géométriques des caractères de \(\mathbb Z/12\mathbb Z\)
s'annulent à moins que \(m-3\) ou \(m+3\) ne soit un multiple de douze. Les trois
amplitudes restantes sont celles indiquées.
\end{proof}

Pour une couronne à colatitude fixe, le support permet uniquement \(m=0\) dans
\(\ell=2\), tandis que dans \(\ell=3\) peuvent subsister
\(m=0,\pm3\). La valeur \(\varepsilon=0\) retrouve le gabarit uniforme
axial ; \(\varepsilon\neq0\) introduit la modulation triadique. Le résultat
réconcilie les deux supports en une enveloppe commune, mais ne fixe pas les amplitudes
ni ne démontre l'exclusivité : le mode \(m=0\) peut également persister pour
\(\ell=3\).

\section{Prédiction d'anisotropie dodécaphasique : paramétrisation préalable, confrontation en aveugle et critère d'incompatibilité observationnelle}
\label{sec:interfaz-cosmologica-c12}

La promotion au rang de prédiction d'anisotropie exige de figer avant d'observer
le ciel :
\begin{enumerate}
 \item l'application \(\mathsf C_{12}^{\rm cyc}\to S^2\) et son équivariance ;
 \item le repère céleste et l'axe, ou ensemble fini d'axes ;
 \item \(\varepsilon,\delta,\phi_0\), la colatitude et la normalisation ;
 \item les \(a_{\ell m}^{\rm HMT}\), leur parité et la statistique
 d'alignement ;
 \item le masque, les avant-plans, le faisceau, le bruit et la covariance ;
 \item la pénalisation pour sélection entre axes ;
 \item la confrontation conjointe en aveugle en température et en polarisation.
\end{enumerate}
La sélection rétrospective de l'axe ou l'ajustement à \(C_\ell^{\rm CMB}\) est une
\textsc{reconnaissance externe}, non une prédiction. L'objet
\(\mathsf C_{12}^{\rm cyc}\) n'est pas \(U_{12}^{\rm sgn}\), une
représentation de \(SU(12)\) ni le spectre angulaire
\(C_\ell^{\rm CMB}\).

\begin{criterion}[Anisotropie dodécaphasique]
\label{crit:anisotropia-dodecafasica}
L'interface est réfutée si le motif scellé est incompatible avec la
confrontation en aveugle ou si l'axe, l'amplitude ou la modulation sont sélectionnés après
l'ouverture des données. Le théorème de support demeure valide même lorsque sa
réalisation cosmologique échoue.
\end{criterion}
